\let\stop\empty
\documentclass{exam-zh}
\usepackage{siunitx}
\usepackage{tkz-euclide}
\usepackage{diagbox}  
\newenvironment{answer}{\par\vspace{0.5em}\noindent\textbf{【答案】}\par}{\par\vspace{1em}}

\newenvironment{solutions}{\par\vspace{0.5em}\noindent\textbf{【解析】}\par}{\par\vspace{1em}}\usepackage{lastpage}
\usepackage{caption}
\examsetup{
  page/size=a4paper,
  paren/show-paren=true,
  paren/show-answer=true,
  fillin/show-answer=false,
  font      = times,
  math-font = xits,
  fillin/no-answer-type=none,
  %solution/show-solution=show-stay,
  %solution/label-indentation=false
  }
%\newcommand{\customcurve}[1][1]{%
  %\tikz[baseline,scale=#1,line width=0.8pt,line cap=round]{%
    %\draw (0.15,0.1) ellipse (0.15 and 0.1);
    %\draw (0,0.1) -- (-0.12 , -0.1);
    %\draw (0,0.1) -- (-0.12 , 0.3);
  %}%
%}
\newcommand{\customcurve}[1][1]{%
  \tikz[baseline,yshift=3,scale=0.05,thick,line width=0.8pt,
        line cap=round, domain=-1:2.832, samples=300]{
    \draw plot (\x,{sqrt(16/(\x+2)^2 - (\x-2)^2)});
    \draw plot (\x,{-sqrt(16/(\x+2)^2 - (\x-2)^2)});
    }
}

\everymath{\displaystyle}

\title{2025 年普通高等学校招生全国统一考试(新课标I卷)}

\subject{数学}



\begin{document}

% \information{
%   姓名\underline{\hspace{6em}},
%   座位号\underline{\hspace{15em}}
% }
% \warning{（在此卷上答题无效）}

\secret

\maketitle
\begin{center}
    本试卷共 \pageref{LastPage} 页，19 题。全卷满分 150 分。考试用时 120 分钟。
\end{center}


\begin{notice}
  \item 答题前，先将自己的姓名、准考证号、考场号、座位号填写在试卷和答题卡上，
    并将准考证号条形码粘贴在答题卡上的指定位置。
  \item 选择题的作答：每小题选出答案后，用 2B 铅笔把答题卡上对应题目的答案标号涂黑。
    写在试卷、草稿纸和答题卡上的非答题区域均无效。
  \item 填空题和解答题的作答：用黑色签字笔直接答在答题卡上对应的答题区域内。
    写在试卷、草稿纸和答题卡上的非答题区域均无效。
  \item 考试结束后，请将本试卷和答题卡一并上交。
\end{notice}



\section{
  选择题：本题共 8 小题，每小题 5 分，共 40 分。
  在每小题给出的四个选项中，只有一项是符合题目要求的。
}

% 1.
\begin{question}
  $(1+5 \iu) \iu$ 的虚部为 \paren
  \begin{choices}
    \item $-1$
    \item 0
    \item 1
    \item 6
  \end{choices}
\end{question}

% 2.
\begin{question}
  设全集 $U=\{x \mid x~ \text{是小于 9 的正整数}\}$，若集合 $A=\{1,3,5\}$ ，则 $\complement_U A$ 中元素个数为 \paren
  \begin{choices}
    \item 0
    \item 3
    \item 5
    \item 8
  \end{choices}
\end{question}

% 3.
\begin{question}
  若双曲线 $C$ 的虚轴长为实轴长的 $\sqrt{7}$ 倍，则 $C$ 的离心率为 \paren
  \begin{choices}
    \item $\sqrt{2}$
    \item $2$
    \item $\sqrt{7}$
    \item $2 \sqrt{2}$
  \end{choices}
\end{question}

% 4.
\begin{question}
  若点 $(a, 0)~(a>0)$ 是函数 $y=2 \tan \left(x-\frac{\uppi}{3}\right)$ 的图象的一个对称中心，则 $a$ 的最小值为 \paren
  \begin{choices}
    \item $30^{\circ}$
    \item $60^{\circ}$
    \item $90^{\circ}$
    \item $135^{\circ}$
  \end{choices}
\end{question}

% 5.
\begin{question}
  设 $f(x)$是定义在$\mathbf{R}$上且周期为2的偶函数, 当$2 \le x \le 3$时, $f(x)=5-2x$, 则$f\left(-\frac{3}{4}\right)=$ \paren
  \begin{choices}
    \item $-\frac{1}{2}$
    \item $-\frac{1}{4}$
    \item $\frac{1}{4}$
    \item $\frac{1}{2}$
  \end{choices}
\end{question}


\newpage
% 6.
\begin{question}
  帆船比赛中，运动员可借助风力计测定风速的大小和方向，测出的结果在航海学中称为视风风速，视风风速对应的向量，是真风风速对应的向量与船行风速对应的向量之和，其中船行风速对应的向量与船速对应的向量大小相等，方向相反. 表\ref{tab1}给出了部分风力等级、名称与风速大小的对应关系. 已知某帆船运动员在某时刻测得的视风风速对应的向量与船速对应的向量如图\ref{fig1}风速的大小和向量的大小相同，单位 $(\mathrm{m} / \mathrm{s})$，则真风为 \paren
  \begin{choices}
    \item 轻风
    \item 微风
    \item 和风
    \item 劲风
  \end{choices}
  
\begin{minipage}[b]{0.4\linewidth}
  \centering
  \begin{tabular}{|c|c|c|}
    \hline
    等级 & 风速大小 $\mathrm{m}/\mathrm{s}$ & 名称 \\
    \hline
    2 & $1.1 \sim* 3.3$ & 轻风 \\
    \hline
    3 & $3.4 \sim* 5.4$ & 微风 \\
    \hline
    4 & $5.5 \sim* 7.9$ & 和风 \\
    \hline
    5 & $8.0 \sim* 10.1$ & 劲风 \\
    \hline
  \end{tabular}
  \captionof{table}{风力等级对照表}\label{tab1}
\end{minipage}
\hspace{2em}
\begin{minipage}[b]{0.4\linewidth}
  \centering
  \begin{tikzpicture}[scale=0.8]
    \tkzInit[xmin=0,xmax=3,ymin=0,ymax=3]
    \tkzGrid
    \tkzDrawX[label=$x$, below right]
    \tkzDrawY[label=$y$, above left]
    \tkzDefPoints{0/0/O,3/3/P,2/0/A,0/2/B}
    \tkzLabelPoint[below left](O){$O$}
    \foreach \i in {1, 2, 3} {
      \draw (\i, 0.1) -- (\i, -0.1) node[below=2pt] {\i};
      \draw (0.1, \i) -- (-0.1, \i) node[left=2pt] {\i};
    }
    \tkzDrawSegment[thick,-latex](A,P)
    \tkzDrawSegment[thick,-latex](P,B)
    \tkzLabelPoint[right](2.5,1.5){船速}
    \tkzLabelPoint[anchor=-40](1.5,2.5){视风速}
  \end{tikzpicture}
  \captionof{figure}{风速矢量图示}\label{fig1}
\end{minipage}
\end{question}


% 7.
\begin{question}
  若圆 $x^2 + (y+2)^2 = r^2~(r > 0)$ 上到直线$y= \sqrt{3}x+2$的距离为1的点有且仅有2个, 则$r$的取值范围是 \paren
  \begin{choices}
    \item $(0,1)$
    \item $(1,3)$
    \item $(3,+\infty)$
    \item $(0,+\infty)$
  \end{choices}
\end{question}

% 8.
\begin{question}
  若实数$x,y,z$ 满足 $2+ \log_2 x = 3 + \log_3 y = 5 + \log_5 z$, 则$x,y,z$ 的大小关系不可能是 \paren
  \begin{choices}
    \item $x > y>z$
    \item $x > z >y$
    \item $y > x>z$
    \item $y > z>x$
  \end{choices}
\end{question}


\section{
  选择题：本题共 3 小题，每小题 6 分，共 18 分。
  在每小题给出的选项中，有多项符合题目要求的。
  全部选对的得 6 分，部分选择的得部分分，有选错的得 0 分。
}
% 9.
\begin{question}
  在正三棱柱 $ABC-A_1B_1C_1$中, $D$为$BC$ 中点, 则 \paren
  \begin{choices}[columns = 2]
    \item $AD \perp A_1C$
    \item $BC \perp \text{平面 } AA_1D$
    \item $AD \parallel A_1B_1$ 
    \item $CC_1 \parallel \text{平面 } A_1AD$
  \end{choices}
\end{question}

% 10.
\begin{question}
  设抛物线 $C: y^2=6 x$ 的焦点为 $F$ ，过 $F$ 的直线交 $C$ 于 $A$、$B$ ，过 $F$ 且垂直于 $A B$ 的直线交准线 $l: x=-\frac{3}{2}$ 于 $E$ ，过点 $A$ 作准线 $l$ 的垂线，垂足为 $D$ ，则 \paren
  \begin{choices}
    \item $|AD| = |AF|$
    \item $|AE| = |AB|$
    \item $|AB| \ge 6$
    \item $|AE| \cdot |BE| \ge 18$
  \end{choices}
\end{question}

% 11.
\begin{question}
  已知$\triangle ABC$ 的面积为$\frac{1}{4}$, 若$\cos 2A + \cos 2B + \cos 2C = 2$,~ $\cos A \cos B \sin C = \frac{1}{4}$, 则 \paren
  \begin{choices}
    \item $\sin C = \sin^2 A + \sin^2 B$
    \item $AB = \sqrt{2}$
    \item $\sin A + \sin B = \frac{\sqrt{6}}{2}$
    \item $AC^2 + BC^2 = 3$
  \end{choices}
\end{question}






\section{填空题：本题共 3 小题，每小题 5 分，共 15 分。}

% 12.
\begin{question}
  若直线$y=2x+5$是曲线$y=\eu^x +x+a$的切线, 则$a=$ \fillin.
\end{question}

% 13.
\begin{question}
  若一个正项等比数列的前 4 项和为 4 ，前 8 项和为 68 ，则该等比数列的公比为 \fillin.
\end{question}

% 14.
\begin{question}
  一个箱子里有 5 个相同的球，分别以 $1 \sim* 5$ 标号，从中有放回地取三次，记至少取出一次的球的个数 $X$，则数学期望 $E(X)=$\fillin.
\end{question}




\section{解答题：本题共 5 小题，共 77 分。解答应写出文字说明、证明过程或者演算步骤。}

%15
\begin{problem}[points = 13]

为研究某疾病与超声波检查结果的关系，从做过超声没检查的人群中随机调查了 1000人，得到如下列联表：
\begin{center}
\begin{tabular}{|c|c|c|c|}
  \hline
  \diagbox[innerleftsep=2cm]{组别}{超声波检查结果}
      & 正常 & 不正常 & 合计 \\
  \hline
  患该疾病   & 20  & 180 & 200 \\
  \hline
  未患该疾病 & 780 & 20  & 800 \\
  \hline
  合计       & 800 & 200 & 1000 \\
  \hline
\end{tabular}
\end{center}
\begin{enumerate}
    \item[(1)] 记超声波检查结果不正常者患该疾病的概率为 $P$ ，求 $P$ 的估计值;
    \item[(2)] 根据小概率值 $\alpha=0.001$ 的独立性检验，分析超声波检查结果是否与患该疾病有关.
  \end{enumerate}

附：$\chi^2=\frac{n(a d-b c)^2}{(a+b)(c+d)(a+c)(b+d)}$,
\begin{center}
\begin{tabular}{|c|c|c|c|}
\hline$P\left(\chi^2 \geq k\right)$ & 0.050 & 0.010 & 0.001 \\
\hline$k$ & 3.841 & 6.635 & 10.828 \\
\hline
\end{tabular}
\end{center}
\end{problem}

%16
\begin{problem}[points = 15]

设数列$\{a_n\}$满足$a_1=3$,~$\frac{a_{n+1}}{n} = \frac{a_n}{n+1} + \frac{1}{n(n+1)}$.
\begin{enumerate}
    \item[(1)] 证明: $\{na_n\}$为等差数列;
    \item[(2)] 设 $f(x)=a_1x+a_2x^2 + \dots + a_m x^m$, 求$f^{\prime}(-2)$.
  \end{enumerate}
\end{problem}


%17
\begin{problem}[points = 15]

如图所示的四棱锥 $P-ABCD$中, $PA \perp$平面 $ABCD$, $BC \parallel AD$, $AB \perp AD$.
  \begin{enumerate}
    \item[(1)] 证明: 平面$PAB \perp$平面$PAD$;
    \item[(2)] 若 $PA=AB=\sqrt{2}$,~ $AD=\sqrt{3}+1$,~ $BC=2$,~ $P, B, C, D$ 在同一个球面上, 设该球面的球心为$O$.
    \begin{enumerate}
        \item[(i)] 证明: $O$在平面$ABCD$上;
        \item[(ii)] 求直线$AC$与直线$PO$所成角的余弦值.
    \end{enumerate}
  \end{enumerate}
  \textfigure{
  \phantom{很长长长长长长长长长长长长长长长长长长长长长长长长长长长长长长长长长长的文本控制图片位置}
  }{
    \includegraphics[width=6cm]{fig/path41.pdf}
  }
\end{problem}

%18
\begin{problem}[points = 17]

设椭圆$C:\frac{x^2}{a^2} + \frac{y^2}{b^2} = 1~(a>b>0)$, 记$A$为椭圆下端点, $B$为右端点, $|AB|=\sqrt{10}$, 且椭圆$C$的离心率为$\frac{2\sqrt{2}}{3}$.
  \begin{enumerate}
    \item[(1)] 求椭圆的标准方程;
    \item[(2)] 设点$P(m,n)$.
    \begin{enumerate}
        \item[(i)] 若$P$不在$y$轴上, 设$R$是射线$AP$上一点, $|AR|\cdot|AP|=3$, 用$m,n$表示点$R$的坐标;
        \item[(ii)] 设直线$OQ$的斜率为$k_1$, 直线$OP$的斜率为$k_2$, 若$k_1=3k_2$, $M$为椭圆上一点, 求$|PM|$的最大值.
    \end{enumerate}
  \end{enumerate}
\end{problem}

%19
\begin{problem}[points = 17]

设函数$f(x) = 5\cos x - \cos 5x$.
  \begin{enumerate}
    \item[(1)] 求$f(x)$在$\left[0, \frac{\uppi}{4}\right]$的最大值;
    \item[(2)] 给定$\theta \in (0,\uppi)$,~ $a$为给定实数, 证明: 存在 $y \in [a-\theta, a+\theta]$, 使得$\cos y \le \cos\theta$;
    \item[(3)] 设$b\in \mathbf{R}$若存在实数$\varphi$, 使得对任意实数$x$都有$5\cos x - \cos(5x+\varphi) \le b$, 求$b$的最小值.
  \end{enumerate}
\end{problem}
%\begin{solution}
  %解答
%\end{solution}


\end{document}